\section{Institutional Background}\label{sec:institutional}

Brazil's federal SME statute (Complementary Law 123/2006, amended 2014) defines micro- and small enterprises by annual gross revenue---R\$360{,}000 (US\$103{,}000) and R\$4{,}800{,}000 (US\$1.37 million) ceilings, respectively---and mandates exclusive SME tenders for items valued at R\$80{,}000 (US\$22{,}900) or less. Public buyer units (PBUs)---federal, state, and municipal administrative entities subject to procurement rules---can avoid the SME-only default only by justifying the departure through a formal report to the audit courts, a costly and scrutinized process. The 2014 amendment converted what had been an option (``may'') into a binding default (``must'') and imposed these \emph{opt-out costs}; state-level adherence to SME-only tendering for eligible items rose from roughly 13\% to nearly 70\% in the subsequent years.\footnote{Adoption-rate figures and Group-65 sub-rates reported throughout this section are computed by the author from the BEC microdata extract used for the structural analysis (Section~\ref{sec:data}), restricted to items satisfying the federal SME-eligibility threshold of R\$80{,}000.}

Since 2005, all public procurement in the state of S\~ao Paulo passes through Bolsa Eletr\^onica de Compras (BEC), a centralized electronic reverse-auction platform in which registered suppliers submit sealed bids (phase 1) followed by an iterative electronic auction (phases 2 onward) in which each participant may progressively lower its offer until the auction closes.\footnote{Brazilian institutional terms used throughout the paper (Preg\~ao, Convite, BEC, CADMAT, CMED, PGE-SP, TCE-SP, Sefaz-SP) are defined in Appendix~\ref{app:glossary}.} BEC handled approximately R\$13 billion (US\$3.7 billion) of procurement in 2019 alone and has cumulatively moved more than R\$105 billion (US\$30 billion) since implementation, spanning 860{,}000 purchase offers and 4.8 million item-level transactions.\footnote{Platform totals are reported by the Secretaria da Fazenda e Planejamento do Estado de São Paulo (Sefaz-SP) and the BEC operating unit at \url{www.bec.sp.gov.br}; the 4.8-million item-level count and the 860{,}000 purchase-offer count are author tallies from the same BEC extract used in the structural analysis.} BEC organizes its catalog into three tiers---group, class, and item---so that ``medical, dental, and hospital equipment and supplies'' form Group 65, within which class 6531 identifies medications subject to price ceilings set by the C\^amara de Regula\c{c}\~ao do Mercado de Medicamentos (CMED).

Group 65 was a singular exception to the 2014 rule. Between August 2014 and February 2018, the state government and the Tribunal de Contas do Estado (TCE-SP) operated under a joint interpretation that medical supplies were strategic goods requiring open competition; during that window no Group-65 bid was restricted to SMEs. In November 2017, the Procuradoria-Geral do Estado (PGE-SP, a separate control agency) issued a legal opinion reversing the exception on isonomy grounds---the constitutional requirement that comparable cases be treated alike. The opinion argued that no substantive feature of medical supplies justified procedural distinctions from other product groups. Effective March 2018, opt-out costs extended to Group 65; SME-only adherence within the group jumped to 43\% on average, below the rate for other groups largely because class-6531 pharmaceuticals are oligopolistic and often fail the three-eligible-SMEs requirement.

What matters for identification is that the policy change was triggered by a legalistic reinterpretation of the isonomy principle at PGE-SP, not by anything happening in the Group-65 procurement market. Public documentation turns up no sign that the decision was a response to price levels, competition patterns, or supplier complaints specific to Group 65 in 2017; the timing was set by the bureaucratic process inside PGE-SP and by when that office chose to communicate the opinion to PBUs. The change is therefore plausibly exogenous to the outcomes studied. This is the institutional basis for Assumption~\ref{a:setaside}: the cutoff shifts admissibility and equilibrium entry selection, not the underlying production technology of the goods themselves. The reduced-form DiD against the 76 never-treated product groups, reported in Appendix~\ref{app:did}, provides one empirical check; the primitive-invariance KS test in Table~\ref{tab:v3_primitive_invariance} provides another.
